\subsection{Finite Sums}

\begin{definition} \label{def:finite-sums-indexed-by-integers} \hlt{Finite Sums Indexed by Integers}\\
Let $a_i\in\C$ be assigned for integers $m\leq i\leq n$. Define $\sum\limits_{i=m}^{n}a_i\defeq a_m+a_{m+1}+\cdots+a_n$.\\
If $m>n$, the sum over the empty index set is $0$.
\end{definition}

\begin{proposition} \label{prop:finite-sum-basic-laws} \hlt{Basic Laws for Finite Sums}\\
Let $a_i,b_i\in\C$ and $c\in\C$ be assigned for $m\leq i\leq n$.
\begin{enumerate}[label=(\roman*)]
\setlength{\itemsep}{0pt}
\item If $m\leq r<n$, then $\sum\limits_{i=m}^{n}a_i=\sum\limits_{i=m}^{r}a_i+\sum\limits_{i=r+1}^{n}a_i$.
\item $\sum\limits_{i=m}^{n}(a_i+b_i)=\sum\limits_{i=m}^{n}a_i+\sum\limits_{i=m}^{n}b_i$ and $\sum\limits_{i=m}^{n}ca_i=c\sum\limits_{i=m}^{n}a_i$.
\item $\abs{\sum\limits_{i=m}^{n}a_i}\leq\sum\limits_{i=m}^{n}\abs{a_i}$.
\item If $a_i,b_i\in\R$ and $a_i\leq b_i$ for every $i$, then $\sum\limits_{i=m}^{n}a_i\leq\sum\limits_{i=m}^{n}b_i$.
\end{enumerate}
\begin{proof}
\begin{enumerate}[label=(\roman*)]
\setlength{\itemsep}{0pt}
\item By Definition \ref{def:finite-sums-indexed-by-integers} and associativity of addition in $\C$,
$\sum\limits_{i=m}^{n}a_i=\sum\limits_{i=m}^{r}a_i+\sum\limits_{i=r+1}^{n}a_i$.
\item This follows from associativity, commutativity, and distributivity in the field $\C$.
\item Induct on $n-m$. The case $n=m$ is immediate.\\
If the claim holds up to $n$, then the triangle inequality gives
\[
\abs{\sum\limits_{i=m}^{n+1}a_i}
\leq \abs{\sum\limits_{i=m}^{n}a_i}+\abs{a_{n+1}}
\leq \sum\limits_{i=m}^{n}\abs{a_i}+\abs{a_{n+1}}
=\sum\limits_{i=m}^{n+1}\abs{a_i}.
\]
\item Induct on $n-m$. The case $n=m$ is the hypothesis $a_m\leq b_m$.\\
If the claim holds up to $n$, then compatibility of the order on $\R$ with addition gives
\[
\sum\limits_{i=m}^{n+1}a_i
=\sum\limits_{i=m}^{n}a_i+a_{n+1}
\leq \sum\limits_{i=m}^{n}b_i+b_{n+1}
=\sum\limits_{i=m}^{n+1}b_i.
\]
\end{enumerate}
\hlt{Structural Commentary}: Finite sums are iterated addition, so their laws are inherited from the scalar field.\\
\hlt{Key Insight}: No convergence hypothesis is needed for a finite sum.
\end{proof}
\end{proposition}

\begin{definition} \label{def:finite-set-summation} \hlt{Summation over a Finite Set}\\
Let $X$ be a finite set with $\#X=n$, and let $f:X\to\C$.\\
Choose any bijection $g:J_n\to X$, where $J_n=\{1,\ldots,n\}$, and set $\sum\limits_{x\in X}f(x)\defeq\sum_{i=1}^{n}f(g(i))$.\\
When $X=\emptyset$, this is the empty sum $0$.\\
\hlt{Key Insight}: Proposition \ref{prop:finite-set-summation-laws} shows that the value is independent of the chosen bijection.
\end{definition}

\begin{proposition} \label{prop:finite-set-summation-laws} \hlt{Finite Set Summation Laws}\\
Let $X,Y$ be finite sets.
\begin{enumerate}[label=(\roman*)]
\setlength{\itemsep}{0pt}
\item \hlt{Well-definedness}: different bijections $J_{\#X}\to X$ give the same value in Definition \ref{def:finite-set-summation}.
\item \hlt{Substitution}: if $h:Y\to X$ is a bijection and $f:X\to\C$, then $\sum\limits_{x\in X}f(x)=\sum\limits_{y\in Y}f(h(y))$.
\item \hlt{Linearity}: for $f,g:X\to\C$ and $c\in\C$, $\sum\limits_{x\in X}(f(x)+g(x))=\sum\limits_{x\in X}f(x)+\sum\limits_{x\in X}g(x)$ and $\sum\limits_{x\in X}cf(x)=c\sum\limits_{x\in X}f(x)$.
\item \hlt{Disjoint unions}: if $X\cap Y=\emptyset$ and $f:X\cup Y\to\C$, then $\sum\limits_{z\in X\cup Y}f(z)=\sum\limits_{x\in X}f(x)+\sum\limits_{y\in Y}f(y)$.
\item \hlt{Finite Fubini}: if $f:X\times Y\to\C$, then $\sum\limits_{x\in X}\sum\limits_{y\in Y}f(x,y)=\sum\limits_{(x,y)\in X\times Y}f(x,y)=\sum\limits_{y\in Y}\sum\limits_{x\in X}f(x,y)$.
\end{enumerate}
\begin{proof}
\begin{enumerate}[label=(\roman*)]
\setlength{\itemsep}{0pt}
\item For $f:X\to\C$, let $n=\#X$ and let $g,k:J_n\to X$ be bijections. Induct on $n$. The case $n=0$ is empty. If $n\geq1$, let $j\in J_n$ satisfy $k(j)=g(n)$. Removing $n$ from $g$ and $j$ from $k$ gives bijections onto $X\setminus\{g(n)\}$; after relabeling by $J_{n-1}$, the induction hypothesis makes the remaining sums equal. Adding the common last term $f(g(n))=f(k(j))$ gives equality of the two sums.
\item Let $g:J_{\#Y}\to Y$ be a bijection. Then $h\circ g:J_{\#Y}\to X$ is a bijection, so $\#Y=\#X$ by Proposition \ref{prop:finite-cardinality-well-defined}.\\
Hence, by (i) and Definition \ref{def:finite-set-summation}, $\sum\limits_{x\in X}f(x)=\sum\limits_{i=1}^{\#Y}f(h(g(i)))=\sum\limits_{y\in Y}f(h(y))$.
\item Choose a bijection $\gamma:J_{\#X}\to X$.\\
Proposition \ref{prop:finite-sum-basic-laws} gives $\sum\limits_{x\in X}(f(x)+g(x)) =\sum\limits_{x\in X}f(x)+\sum\limits_{x\in X}g(x)$ and $\sum\limits_{x\in X}cf(x)=c\sum\limits_{x\in X}f(x)$
\item Let $\alpha:J_{\#X}\to X$ and $\beta:J_{\#Y}\to Y$ be bijections, and enumerate $X\cup Y$ by listing $\alpha$ first and then $\beta$. Since $X\cap Y=\emptyset$, this is a bijection, and Proposition \ref{prop:finite-sum-basic-laws} gives $\sum\limits_{z\in X\cup Y}f(z)=\sum\limits_{x\in X}f(x)+\sum\limits_{y\in Y}f(y)$.
\item Let $\alpha:J_{\#X}\to X$ and $\beta:J_{\#Y}\to Y$ be bijections.\\
The maps $\ell((i-1)\#Y+j)=(\alpha(i),\beta(j))$ and $r((j-1)\#X+i)=(\alpha(i),\beta(j))$ for $i\in J_{\#X}$ and $j\in J_{\#Y}$ give the two lexicographic bijections $J_{(\#X)(\#Y)}\to X\times Y$.\\
Applying Definition \ref{def:finite-set-summation}, Proposition \ref{prop:finite-sum-basic-laws}, and (i) yields $\sum\limits_{x\in X}\sum\limits_{y\in Y}f(x,y) =\sum\limits_{(x,y)\in X\times Y}f(x,y) =\sum\limits_{y\in Y}\sum\limits_{x\in X}f(x,y)$.
\end{enumerate}
\hlt{Structural Commentary}: A finite set sum is an order-free finite sum.\\
\hlt{Key Insight}: Bijections are bookkeeping devices; the value does not depend on the bookkeeping order.
\end{proof}
\end{proposition}
